% Scientific Reports front matter.
% Abstract: <= 200 words, unstructured, a general introduction and a brief
% non-technical summary of the main results and their implications; no
% citations, no figures. Keywords: <= 6. No plain-language summary in this
% journal. The Communications Medicine abstract is 01_abstract_CM.tex.
\begin{abstract}
\noindent
Medical machine learning increasingly asks one kind of data to report on
something it does not directly measure: a tissue section on a tumour's
mutations, a brain scan on its molecular markers, an electrocardiogram
on a patient's physiology. Such predictors are judged by one
cohort-average accuracy per target, a number that cannot say whether
the target is reachable at all, or what to change if not. Here we report per-target \emph{translatability} instead.
Pairing an expert vocabulary of the source data with an unconstrained
learned encoder gives two sources of evidence about each target, and
the four ways two sources can share information define three
regimes---encoder saturation, complementarity and
out-of-ontology---each with a different remedy. Regimes are read from
the direction of contrasts between four retrained models, adjudicated
by a hierarchical bootstrap, and issued as a Conclusive/Inconclusive
verdict at the level at which it was established. On 444 lung
adenocarcinoma patients, the scarcity typical of public cohorts, the
instrument resolves one contrast, declines to assign any regime, and
reports the cohort size that could. The
same protocol runs unchanged on glioma MRI, declining MGMT at full
power, on the electrocardiogram, and across cohorts, where it turns
Inconclusive rather than confidently wrong.
\end{abstract}

\paragraph{Keywords.}
Cross-domain prediction; partial information decomposition; model
evaluation; interpretability; multiple-instance learning; computational
pathology.
